\section{Vehicles}

\subsection{AX3 Razorshark Fighter}

For Fire Warriors on the front line, few sights are more comforting than the arrival of a Razorshark Fighter bursting from the cloud ceiling and diving toward enemy formations. Its chassis has many elements in common with that of the Sun Shark Bomber, except that the Razorshark forgoes interceptor Drones in favour of additional manoeuvre engines, and exchanges the Pulse Bomb Generator for a formidable quad ion turret. The latter is a terrible weapon against light vehicles, which it promptly turns into twisted wrecks. At the cost of slight instability, the quad ion turret can also fire in overload mode, and thus affect an area sufficient to consume entire squads in a single blast.

The Razorshark has the \textbf{Integral Armour}, \textbf{Interceptor}, \textbf{Jink (5+)}, and \textbf{Flyer} abilities. Its Ion Cannons have the \textbf{Turret} ability.

\begin{unitstatstable}
\SetCell[r=2]{c,m} -- &
\SetCell[r=2]{c,m} 3+ &
\SetCell[r=2]{c,m} +3 &
\SetCell[r=2]{c,m} 6 &
Quad Ion Cannons &
30 cm &
2 &
3+ &
$-1$
\\
& & & &
Burst Cannons &
20 cm &
2 &
4+ &
0
\end{unitstatstable}

\unitimage{razorshark.png}



\unitdivider

\subsection{AX39 Sun Shark Bomber}

Their combat doctrines resting on mobility, Hunter Cadres do not employ artillery in the conventional sense. Instead, in collaboration with the Air Caste, they have developed the Sun Shark Bomber, an elegant atmospheric aircraft capable of plunging from the skies to strike ground targets with its devastating armament. Sun Sharks are therefore the alpha predators of the skies in which they operate. The mere sight of a formation of such bombers, arriving in graceful flight over the battlefield while leaving behind a carpet of bluish explosions, is enough to plunge the enemies of the T'au Empire into indescribable terror.

It has the \textbf{Integral Armour} and \textbf{Flyer} abilities. Its bombs have the \textbf{Bombardment (1)} and \textbf{Template (7.5 cm)} abilities. Its missiles have the \textbf{Smart Missiles} ability.

\begin{unitstatstable}
\SetCell[r=2]{c,m} -- &
\SetCell[r=2]{c,m} 3+ &
\SetCell[r=2]{c,m} +2 &
\SetCell[r=2]{c,m} 6 &
Smart Missiles &
30 cm &
3 &
4+ &
$-1$
\\
& & & &
Pulse Bombs &
Bomb &
Template &
4+ &
$-2$
\end{unitstatstable}

\unitimage{sun-shark.png}



\unitdivider

\subsection{Barracuda}

The Barracuda air-superiority fighter is the most common atmospheric aircraft of the T'au. Faster than the Marauder Bomber of the Imperial Navy, but without the top speed of the Thunderbolt or Lightning Fighter, it compensates through its sophisticated electronic systems and the natural talents of its pilots. In aerial combat, Thunderbolts and Barracudas are very close, the Thunderbolt being the faster and the Barracuda the more manoeuvrable. Imperial pilots often have combat experience that the T'au Air Caste struggles to match.

The Barracuda has the \textbf{Integral Armour}, \textbf{Jink (5+)}, \textbf{Interceptor}, and \textbf{Flyer} abilities. Its Burst Cannons have the \textbf{Turret} ability.

\begin{unitstatstable}
\SetCell[r=2]{c,m} -- &
\SetCell[r=2]{c,m} 3+ &
\SetCell[r=2]{c,m} +3 &
\SetCell[r=2]{c,m} 6 &
Ion Cannons &
45 cm &
1 &
4+ &
$-2$
\\
& & & &
Burst Cannons &
20 cm &
2 &
4+ &
0
\end{unitstatstable}

\unitimage{barracuda.png}



\unitdivider

\subsection{DX6 Remora Fighter Drone}

The DX6 Remora Fighter is a stealth combat aircraft commanded by a drone armed with two twin burst cannons, perfect for striking lightly armoured units.

The Remora has the \textbf{Integral Armour}, \textbf{Camouflage}, \textbf{Markerlight (45 cm)}, \textbf{Flyer}, and \textbf{Drone} abilities.

\begin{unitstatstable}
30 cm & 5+ & $-1$ & -- & Twin Burst Cannons & 20 cm & 2 & 4+ & $-1$
\end{unitstatstable}

\unitimage{remora.png}



\unitdivider

\subsection{Devilfish}

The armoured Devilfish troop transport is the workhorse of all T'au ground forces, and provides Fire Caste infantry with precious mobility as well as a valuable fire-support unit in all circumstances. A Devilfish can carry up to twelve Fire Caste warriors with their full kit, and take them into combat in relative safety, then provide mobile support once they have disembarked. Highly mechanised cadres rely on Devilfish to move their troops quickly in order to redeploy their firing line to the most crucial place and support Hammerhead tanks.

The Devilfish has the \textbf{Antigrav}, \textbf{Attached Transport}, and \textbf{Transport (2)} abilities.

\begin{unitstatstable}
25 cm & 3+ & +0 & Attached & Burst Cannon & 30 cm & 1 & 4+ & 0
\end{unitstatstable}

\unitimage{devilfish.png}



\unitdivider

\subsection{Interceptor Drones}

Interceptor Drones are flying drones equipped with heavy anti-aircraft weapons.

Interceptor Drones have the \textbf{Integral Armour}, \textbf{Drone}, \textbf{Attached Character}, \textbf{Shield Drone}, and \textbf{Flyer} abilities. They may only be attached to Flyers.

\begin{unitstatstable}
-- & -- & $-1$ & Attached & Ion Cannons & 20 cm & 1 & 4+ & 0
\end{unitstatstable}

\unitimage{interceptor-drones.png}



\unitdivider

\subsection{Sentry Drone}

A Sentry Drone Turret is a group of automatic weapon systems linked together to form a defensive perimeter. Sentry Drone Turrets are deployed from orbit or from very high-altitude vessels and have a small jump reactor to reach their position. Once in place, the turrets remain stationary until they are recovered.

Sentry Drone Turrets have the \textbf{Autonomous Defence}, \textbf{Integral Armour}, \textbf{Drone}, and \textbf{Deep Strike (2)} abilities. Their Burst Cannons have the \textbf{Turret} ability.

\begin{unitstatstable}
0 cm & 5+ & $-2$ & -- & Burst Cannons & 30 cm & 2 & 4+ & 0
\end{unitstatstable}

\unitimage{sentry-drone.png}



\unitdivider

\subsection{Hammerhead}

The T'au main battle tank evokes contained violence and bears the name ``Hammerhead''. Like a starving predator on the prowl, all enemies who have had dealings with it have quickly learned to fear its firepower. First encountered by the Imperium during the Damocles Crusade, the Hammerhead has since been the basic combat vehicle of T'au Cadres, although alternative armaments have since been successfully identified. Hammerhead tanks are commonly encountered in small numbers, to support Fire Warrior teams.

Hammerheads have the \textbf{Antigrav} ability. Their main weapon has the \textbf{Turret} ability; in addition the Railgun of the Mk II variant has the \textbf{Damage (+1)} ability, and the Missile Rack of the Mk III uses \textbf{Artillery} and \textbf{Reduces Cover ($-1$)}. Finally the missiles of the Mk II, III, and VI variants have the \textbf{Smart Missiles} ability.

\subsubsection{Hammerhead Mk I}

\begin{unitstatstable}
\SetCell[r=2]{c,m} 20 cm &
\SetCell[r=2]{c,m} 3+ &
\SetCell[r=2]{c,m} +0 &
\SetCell[r=2]{c,m} 6 &
Ion Stream Cannons &
30 cm &
2 &
3+ &
$-3$
\\
& & & &
Burst Cannons &
30 cm &
2 &
4+ &
0
\end{unitstatstable}

\subsubsection{Hammerhead Mk II}

\begin{unitstatstable}
\SetCell[r=2]{c,m} 20 cm &
\SetCell[r=2]{c,m} 3+ &
\SetCell[r=2]{c,m} +0 &
\SetCell[r=2]{c,m} 6 &
Railgun &
75 cm &
1 &
4+ &
$-3$
\\
& & & &
Smart Missiles &
45 cm &
1 &
4+ &
$-1$
\end{unitstatstable}

\subsubsection{Hammerhead Mk III}

\begin{unitstatstable}
\SetCell[r=2]{c,m} 20 cm &
\SetCell[r=2]{c,m} 3+ &
\SetCell[r=2]{c,m} +0 &
\SetCell[r=2]{c,m} 6 &
Missile Rack &
75 cm &
2 &
3+ &
$-1$
\\
& & & &
Smart Missiles &
45 cm &
1 &
4+ &
$-1$
\end{unitstatstable}

\subsubsection{Hammerhead Mk IV}

\begin{unitstatstable}
\SetCell[r=2]{c,m} 20 cm &
\SetCell[r=2]{c,m} 3+ &
\SetCell[r=2]{c,m} +0 &
\SetCell[r=2]{c,m} 6 &
Heavy Burst Cannons &
45 cm &
4 &
3+ &
0
\\
& & & &
Burst Cannons &
30 cm &
1 &
4+ &
0
\end{unitstatstable}

\subsubsection{Hammerhead Mk V}

\begin{unitstatstable}
\SetCell[r=2]{c,m} 20 cm &
\SetCell[r=2]{c,m} 3+ &
\SetCell[r=2]{c,m} +0 &
\SetCell[r=2]{c,m} 6 &
Ion Cannon &
75 cm &
2 &
4+ &
$-2$
\\
& & & &
Burst Cannons &
30 cm &
2 &
4+ &
0
\end{unitstatstable}

\subsubsection{Hammerhead Mk VI}

\begin{unitstatstable}
\SetCell[r=2]{c,m} 20 cm &
\SetCell[r=2]{c,m} 3+ &
\SetCell[r=2]{c,m} +0 &
\SetCell[r=2]{c,m} 6 &
Twin Plasma Rifles &
45 cm &
2 &
3+ &
$-2$
\\
& & & &
Smart Missiles &
45 cm &
1 &
4+ &
$-1$
\end{unitstatstable}

\unitimage{hammerhead.png}



\unitdivider

\subsection{Sky Ray}

The Sky Ray missile tank is a specialised variant of the Hammerhead whose turret cannon is replaced by a rack of formidable guided missiles. Originally designed to neutralise enemy aircraft, the Sky Ray has also proved itself against ground targets. In the field, the Sky Ray is an antigrav tank that floats at the rear of the front, sweeping the horizon with its two Markerlights in search of a ground or aerial target to lock.

The Sky Ray has the \textbf{Antigrav} and \textbf{Markerlight (75 cm)} abilities. Its weapon has the \textbf{AA} and \textbf{Seeker Missiles} abilities.

\begin{unitstatstable}
20 cm & 3+ & +0 & 6 & Seeker Missiles & 75 cm & 2 & 4+ & $-2$
\end{unitstatstable}

\unitimage{sky-ray.png}



\unitdivider

\subsection{Stingray}

After the Damocles Gulf Crusade, the T'au encounter with old Imperial forces led them to consider developing a weapon that could fire outside and beyond lines of sight to break the waves of Orks unleashed in the sector after the Imperium/T'au ceasefire. A large unique warhead, similar to that used by Earthshaker Cannons, was first envisaged, but the T'au quickly realised it would be less effective than a modified seeker missile, with a large number of smart missiles loaded in the hold. The result was first a tank based on the Sky Ray chassis. Early field trials were promising, but revealed shortcomings in the vehicles used for the test. The first of these was the size of the new missile, which made the tanks extremely dependent on resupply during prolonged battles.

The Stingray has the \textbf{Antigrav} ability. The Submunition Missiles have the \textbf{Seeker Missiles} ability, use a \textbf{Template (7.5 cm)} that \textbf{Reduces Cover ($-1$)}, and have the \textbf{Battery} ability; the Smart Missiles have the \textbf{Smart Missiles} ability.

\begin{unitstatstable}
\SetCell[r=2]{c,m} 20 cm &
\SetCell[r=2]{c,m} 3+ &
\SetCell[r=2]{c,m} +0 &
\SetCell[r=2]{c,m} 6 &
Submunition Missiles &
75 cm &
Template &
3+ &
0
\\
& & & &
Smart Missiles &
45 cm &
1 &
4+ &
$-1$
\end{unitstatstable}

\unitimage{stingray.png}


